% =====================================================================
%  CAPITULO 3 - COMPLEMENTO: LEITURA TOPOLOGICA E NOS EQUIPOTENCIAIS
% =====================================================================
\subsection{Leitura topológica e nós equipotenciais}

\begin{conceito}[Antes de calcular, identifique os nós]
Em redes desenhadas de forma não convencional, a geometria pode enganar. Fios
ideais contínuos pertencem ao mesmo nó, mesmo quando percorrem caminhos longos
ou cruzam o desenho. Um cruzamento só representa conexão elétrica quando houver
indicação explícita de junção; quando um fio ``salta'' sobre o outro, não há
contato. Depois de identificar os nós equipotenciais, redes aparentemente
complexas podem se reduzir a associações simples.
\end{conceito}

\blocoaprofundamento

\begin{questao}[3]{}
No circuito da figura, os dois fios diagonais \textbf{se cruzam sem contato
elétrico} no ponto central. Determine a resistência equivalente vista entre os
terminais $a$ e $b$.

\circuitoq{%
\begin{circuitikz}[scale=0.92,transform shape,line width=0.8pt]
  \coordinate (A) at (0,0);
  \coordinate (B) at (5,0);
  \coordinate (C) at (0,4);
  \coordinate (D) at (5,4);

  \draw (C) to[R,l_={$\SI{8}{\ohm}$}] (A);
  \draw (C) to[R,l={$\SI{6}{\ohm}$}] (D);
  \draw (D) to[R,l={$\SI{12}{\ohm}$}] (B);
  \draw (A) to[R,l_={$\SI{4}{\ohm}$}] (B);

  % Primeiro fio diagonal; o segundo e desenhado por cima com uma faixa branca,
  % deixando claro que existe cruzamento, mas nao existe juncao eletrica.
  \draw (C) -- (B);
  \draw[preaction={draw=white,line width=4.2pt}] (A) -- (D);

  \fill (A) circle (1.8pt);
  \fill (B) circle (1.8pt);
  \fill (C) circle (1.8pt);
  \fill (D) circle (1.8pt);

  \draw (A) -- ++(0,-0.75) node[below,font=\normalsize] {$a$};
  \draw (B) -- ++(0,-0.75) node[below,font=\normalsize] {$b$};
\end{circuitikz}}

\resposta{$R_{eq}=\SI{1.6}{\ohm}$}
\resolucao{O passo decisivo é reconhecer os nós. O fio diagonal que liga o
canto superior esquerdo ao canto inferior direito torna esses dois pontos o
\textbf{mesmo nó}, isto é, o nó de $b$. Da mesma forma, o outro fio diagonal
liga o canto inferior esquerdo ao canto superior direito, tornando-os o
\textbf{mesmo nó}, isto é, o nó de $a$.

Como o cruzamento central não é uma junção, os dois nós permanecem distintos.
Assim, os quatro resistores de $\SI{8}{\ohm}$, $\SI{6}{\ohm}$,
$\SI{12}{\ohm}$ e $\SI{4}{\ohm}$ estão todos conectados entre o mesmo par de
nós $a$--$b$: são quatro resistores em paralelo.
\[
\frac{1}{R_{eq}}
 =\frac{1}{8}+\frac{1}{6}+\frac{1}{12}+\frac{1}{4}
 =\frac{3+4+2+6}{24}
 =\frac{15}{24}=\frac{5}{8}.
\]
Logo,
\[
R_{eq}=\frac{8}{5}=\SI{1.6}{\ohm}.
\]
O desenho parece conter uma rede cruzada, mas a leitura topológica transforma o
problema em um paralelo simples.}
\end{questao}

\begin{questao}[3]{}
No circuito da figura, o aro externo é um \textbf{fio ideal contínuo}. Os oito
resistores radiais têm resistência de $\SI{4}{\ohm}$ cada. Determine a corrente
$i$ fornecida pela fonte ideal de $\SI{30}{\volt}$.

\circuitoq{%
\begin{circuitikz}[scale=0.88,transform shape,line width=0.8pt]
  \coordinate (O) at (0,0);
  \coordinate (P) at ({2.35*cos(140)},{2.35*sin(140)});
  \coordinate (N) at ({2.35*cos(220)},{2.35*sin(220)});

  % Aro externo: um unico no eletrico.
  \draw (O) circle (2.35);

  % Oito resistores radiais iguais.
  \foreach \ang in {0,45,90,135,180,225,270,315}{
    \draw (O) to[R] ({2.35*cos(\ang)},{2.35*sin(\ang)});
  }
  \fill (O) circle (1.8pt);
  \node[fill=white,inner sep=1.5pt,font=\footnotesize\sffamily] at (0,-2.85)
       {8 resistores radiais: $\SI{4}{\ohm}$ cada};

  % Fonte e resistor serie de 2 ohms.
  \draw (-4.15,-1.15) to[battery1,l_={$\SI{30}{\volt}$}] (-4.15,1.15)
        to[R,l={$\SI{2}{\ohm}$}] (P);
  \draw (-4.15,-1.15) -- (N);
  \draw[->,Icol,line width=1pt] (-4.55,-0.55) -- (-4.55,0.55)
       node[midway,left,font=\footnotesize] {$i$};
\end{circuitikz}}

\resposta{$i=\SI{15}{\ampere}$}
\resolucao{Todo o aro externo é feito de fio ideal, portanto todos os seus
pontos pertencem ao \textbf{mesmo nó} e têm o mesmo potencial, que chamaremos
de $V_A$. O centro da roda é um segundo nó, de potencial $V_C$.

Os oito resistores de $\SI{4}{\ohm}$ ligam o nó central $C$ ao mesmo nó externo
$A$. Entretanto, $C$ é um nó flutuante: ele não possui nenhuma outra conexão.
Aplicando a LKC ao centro,
\[
8\,\frac{V_C-V_A}{4}=0
\qquad\Longrightarrow\qquad
V_C=V_A.
\]
Logo, a tensão em cada resistor radial é nula e nenhum deles conduz corrente.
Eles não alteram a resistência vista pela fonte.

O único elemento resistivo efetivamente no caminho fechado da fonte é o
resistor de $\SI{2}{\ohm}$. Assim,
\[
i=\frac{30}{2}=\SI{15}{\ampere}.
\]
A chave do problema é perceber que o aro externo curto-circuita todos os pontos
da circunferência em um único nó; calcular o paralelo dos oito resistores sem
observar que o nó central está flutuante levaria a uma interpretação incorreta.}
\end{questao}

\gabarito[3.5 Leitura topológica e nós equipotenciais]
